\section{Modeling Verifier-Based Synthetic Retraining: the Linear Regression Case}\label{sec:linear}

In this section, we formalize our model of iterative retraining with verified synthetic data, coined \emph{verifier-based synthetic retraining} for convenience. Following recent works in this space \citep{dohmatob2024model,gerstgrasser2024model,garg2025preventing,zhusynthesize}, we focus on the foundational linear regression setting where the objective is to estimate a high-dimensional coefficient vector $\theta^\star$ in the following linear model
\[
y = x^{\top} \theta^\star + \xi,
\]
where $\xi \sim \mathcal{N}(0, \sigma^2)$, $x \in \mathbb{R}^p$, and $\theta^\star \in \mathbb{R}^p$ is the unknown parameter of interest. We use the standard Mean Squared Error (MSE), i.e., $\text{MSE}(\hat{\theta}) = \E_{\xi} \|\hat{\theta} - \theta^\star\|^2$, to evaluate estimators.

\paragraph{Modeling the verifier and data filtering rule.}  Suppose we have access to a verifier that possesses prior knowledge of $\theta^\star$, modeled by a knowledge set. Specifically,  the verifier's knowledge is described by a spherical ball:
\[
B_r(\theta_c) := \bigl\{\, \theta \in \mathbb{R}^p : \|\theta - \theta_c\| \leq r \,\bigr\},
\]
with fixed center $\theta_c$ and radius $r$.  We assume this knowledge set indeed contains the true parameter, i.e.,  $\theta^\star \in B_r(\theta_c)$, but  the true parameter $\theta^\star$ is unknown.  The verifier does not reveal $\theta_c$ or $r$ directly (see modeling motivations below). Instead, it  only provides binary feedback indicating whether a given (real or synthetic) data point $(x_i, y_i)$ is consistent with the knowledge $\theta^\star \in B_r(\theta_c)$ or not. Specifically, the verifier outputs \emph{Yes} if
\begin{align}
|y_i - x_i^{\top} \theta_c| \leq r \|x_i\| + \sigma_c, \label{equ:linear:verify_condition}
\end{align}
and \emph{No} otherwise. Here $\sigma_c$ is a  constant related to the verifier's capability.  This \emph{filtering rule} is motivated by the following bound on expected errors: 
$
\E\bigl[\,|y_i - x_i^{\top}\theta_c|\,\bigr] = \E\bigl[\,|x_i^{\top}(\theta^\star - \theta_c) + \xi_i|\,\bigr] \leq r\|x_i\| + \E|\xi_i| = r\|x_i\| + \sqrt{\tfrac{2}{\pi}}\sigma.
$
Since the true $\sigma$ might be unknown in practice, $\sigma_c$ serves as an estimate of the true $\sigma$. %\hf{Do our theory require $\sigma_c = \sigma$ in order to work? }

We refer to $\Delta = \|\theta^\star - \theta_c\|$ as the \emph{bias} of the verifier, whereas  $r$ captures the \emph{selectivity} of the verifier --  the smaller $r$ is, less likely the verifier  accepts a data point $(x_i, y_i)$.  The verifier only needs to  provide Yes/No answers based on the above selection rule in   \eqref{equ:linear:verify_condition}, but does not  needs    to know the parameter  $\theta_c, r$ of the knowledge set.   The motivation of this modeling primarily comes from practice, as explained below. 

  
\paragraph{Motivation of binary feedback from verifiers.}    
We adopt the  binary feedback  from verifiers mainly for practical reasons.   
In practice, eliciting simple yes/no feedback is far less noisy and more cost-effective than asking verifiers to directly specify $\theta_c$ or $r$. Indeed, in real applications verifiers may not even know these quantities explicitly, which would correspond to model parameters if the verifier is a stronger teacher model or how the human reasons if the verifier is a human.  This model choice   is also aligned with the  widely adopted comparison-based feedback in reinforcement learning from human feedback (RLHF) \citep{ouyang2022training}. Such binary feedback has become a standard approach in preference alignment for large language models, where LLM raters and human evaluators provide pairwise or accept/reject judgments that effectively guide learning at scale \citep{wettigqurating}. Although simple, our theory and empirical evaluations both show that this \emph{single bit} of information for each sample can successfully be injected into the retraining process to improve models. 


\paragraph{Synthetic Retraining with Verifier-based Filtering}

We begin with an initial set of real data $(X^0, Y^0)$, where $X^0 \in \mathbb{R}^{n_0 \times p}$ and $Y^0 \in \mathbb{R}^{n_0}$.  
The initial estimator $\hat{\theta}^0$ is obtained via Ordinary Least Squares (OLS) \footnote{For ease of presentation, we assume Rank$(X^0)=p$. If Rank$(X^0)<p$, all our results are equally applicable by working in the subspace of $X^0$. }
$
\hat{\theta}^0 = ({X^{0}}^\top X^0)^{-1} {X^0}^\top Y^0. \label{equ:linear:theta0}
$
We then proceed with iterative synthetic retraining via the \emph{generate--verify--retrain} procedure outlined in Figure~\ref{fig:pipeline}, the rigorous retraining Scheme~\ref{scheme:linear_retraining} and Algorithm~\ref{alg:linear_retraining} are provided in Appendix~\ref{app:linear}.

\begin{figure}[h]
\centering
\begin{tikzpicture}[
  node distance=14mm, >=Stealth,
  midbox/.style={rectangle, rounded corners, draw, align=center,
                 minimum height=8mm, text width=20mm, inner sep=2mm, font=\footnotesize},
  endbox/.style={rectangle, rounded corners, draw, align=center,
                 minimum height=8mm, text width=10mm, inner sep=2mm, font=\footnotesize}
]
\footnotesize
\node[endbox] (theta) {$\hat{\theta}^k$};
\node[midbox, right=of theta] (xy) {$(X^{k+1},\, Y^{k+1})$};
\node[midbox, right=of xy] (xyv) {$({X^{k+1}}',\, {Y^{k+1}}')$};
\node[endbox, right=of xyv] (theta2) {$\hat{\theta}^{k+1}$};

\draw[->] (theta) -- node[midway, above] {Generate} (xy);
\draw[->] (xy) -- node[midway, above] {Verify} (xyv);
\draw[->] (xyv) -- node[midway, above] {Retrain} (theta2);
\end{tikzpicture}
\caption{Generate-Verify-Retrain pipeline.}
\label{fig:pipeline}
\end{figure}

Since learning proceeds through the conditional $Y^k \mid X^k$, synthetic retraining requires specifying the covariate design $X^k$; labels $Y^k$ are then generated conditionally via the model under verifier constraints.
In principle, one could construct $X^k$ arbitrarily; however, for mathematical clarity, below we describe a targeted though arguably natural design. 
In particular, we choose to align the synthetic covariates with a fixed orthonormal set $\{v_1,\dots,v_p\}$ and construct $X^k$ in a block-structured form by repeating each $v_j^\top$ as rows:
\label{equ:linear:block_design}
\begin{align}
X^k = (\,
\underbrace{v_1, \ldots}_{\text{copies of } v_1},\;
\underbrace{v_2, \ldots}_{\text{copies of } v_2},\;
\ldots,\;
\underbrace{v_p, \ldots}_{\text{copies of } v_p}
\,)^\top.
\end{align}
After verifier filtering, each orthogonal direction $v_j$ retains exactly $n_k$ samples with 
$
n_0 \;\leq\; pn_1 \;\leq\; pn_2 \;\leq\; \cdots.
$

Notably, the estimation of parameter $\hat{\theta}^{k+1}$ using only synthetic data from the model with $\hat{\theta}^{k}$, though with filtering, leads to a Markovian transition   $\hat{\theta}^k \mapsto \hat{\theta}^{k+1}$. The above block design essentially helps ``diagonalize''  the transition operator $\hat{\theta}^k \mapsto \hat{\theta}^{k+1}$.  
The conceptual benefit of this covariance design choice is that we remove the rotational variability that arbitrary designs would introduce across iterations and decouple the dynamics along orthogonal directions.
In practice,  this design mirrors curating data along approximately orthogonal latent spaces or topics (e.g., topical axes like politics, sports, mathematics).
However, the choice of covariant $X^k$ is not unique:   alternatives (e.g., canonical basis, isotropic random directions) can yield similar qualitative conclusions, with potentially different constants or rates.
We expect our theoretical insights to generalize to any reasonable design of the covariant $X^k$, though the rigorous proofs may be less tractable for some designs.



